% Recovered verbatim from chapters_31_35_with_part/chapters_31_35_with_part.tex
% Canonical chapter 33
\chapter{ER-Network Scaling in the Cosmic Web}

\section{Cosmic Adjacency Graph}

On large scales, galaxies, clusters, and filaments form a network that can be idealised as a cosmic graph.
In MK--43 Ultra, this graph is a macroscopic imprint of the underlying ER network, with an effective adjacency matrix
\[
  \mathcal{W}_{IJ} \sim \langle W_{ij} \rangle_{\text{coarse-grained}},
\]
where $I,J$ label coarse-grained regions (e.g.\ halos, filaments, voids).

\section{Scaling Laws and Connectivity}

We may characterise the cosmic adjacency graph by:
\begin{itemize}
  \item degree distributions of $\mathcal{W}_{IJ}$,
  \item clustering coefficients,
  \item percolation thresholds for large-scale connectivity.
\end{itemize}
MK--43 Ultra predicts that these statistics are not arbitrary but reflect the ER-foam dynamics and chromatic evolution across cosmological cycles.

\section{Voids, Filaments, and Superclusters}

Regions of low ER connectivity correspond to cosmic voids, while strongly connected regions trace filaments and superclusters.
The pattern of connectivity and its evolution may lead to testable predictions about the distribution and dynamics of the cosmic web.

